\documentclass[12pt,a4paper]{article}

\usepackage{cmap}
\usepackage[T2A]{fontenc}
\usepackage[utf8]{inputenc}
\usepackage[russian]{babel}
\usepackage{amsmath,amssymb}
\usepackage[a4paper,margin=2cm]{geometry}

\setlength{\parindent}{0pt}
\setlength{\parskip}{0.45em}
\sloppy

\begin{document}

\begin{center}
{\Large\bfseries Билет 6: вопросы для проверки знаний}
\end{center}

\noindent\fbox{%
  \begin{minipage}{\dimexpr\linewidth-2\fboxsep-2\fboxrule\relax}
  \normalfont
Определенный интеграл Римана. Верхние и нижние суммы Дарбу. Свойства сумм
Дарбу. Критерии интегрируемости. Интегрируемость непрерывной функции,
монотонной функции, ограниченной функции с конечным числом точек разрыва.
Аддитивность интеграла по отрезкам, линейность интеграла, интегрируемость
произведения функций, интегрируемость модуля интегрируемой функции,
интегрирование неравенств, теорема о среднем. Свойства интеграла с переменным
верхним пределом: непрерывность, дифференцируемость. Формула
Ньютона--Лейбница. Замена переменных и интегрирование по частям в определенном
интеграле.
  \end{minipage}%
}

\section*{Вопросы на понимание}

\begin{enumerate}
\item Чем верхние и нижние суммы Дарбу отличаются от интегральных сумм Римана с выбранными точками?
\item Как по разности \(S(f;T)-s(f;T)\) понять, что график функции на мелких отрезках почти не ``скачет''?
\item Почему при измельчении разбиения нижняя сумма не может уменьшиться, а верхняя не может увеличиться?
\item Почему ограниченность функции является необходимым условием интегрируемости по Риману на отрезке?
\item Чем критерий \(\Delta(f;T)\to0\) удобнее простой проверки одной последовательности разбиений?
\item Почему непрерывность на всем отрезке гарантирует интегрируемость, а наличие конечного числа разрывов ей не мешает?
\item Почему монотонная функция может быть интегрируемой, даже если у нее есть скачки?
\item Может ли изменение значений функции в конечном числе точек изменить ее интегрируемость или значение интеграла? Почему?
\item Почему из интегрируемости \(f\) следует интегрируемость \(|f|\), но обратное утверждение в общем случае неверно?
\item Почему интегрируемость произведения \(fg\) не следует только из линейности интеграла?
\item Как интерпретировать теорему о среднем для интеграла при \(g\equiv1\), и почему для общей версии важно, чтобы \(g\) не меняла знак?
\item Почему интеграл с переменным верхним пределом всегда непрерывен для интегрируемой функции, но равенство \(F'(x)=f(x)\) требует непрерывности \(f\) в соответствующей точке?
\end{enumerate}

\section*{Источники для сверки}

\texttt{target/answer\_question\_06.tex};
\texttt{IvGE\_1.pdf}, стр. 207--227.

\end{document}
